% TeXlate 评测体系与指标演化全景报告
% 编译: latexmk -xelatex report.tex  (或 xelatex ×2)
\documentclass[11pt,a4paper]{ctexart}

\usepackage{geometry}
\geometry{margin=2.4cm,top=2.6cm,bottom=2.6cm}
\usepackage{pgfplots}
\pgfplotsset{compat=1.18}
\usepgfplotslibrary{dateplot,groupplots,statistics}
\usepackage{tikz}
\usetikzlibrary{positioning,arrows.meta,shapes.geometric,fit,backgrounds,calc}
\usepackage{booktabs}
\usepackage{array}
\usepackage{multirow}
\usepackage{caption}
\captionsetup{font=small,labelfont=bf}
\usepackage{xcolor}
\usepackage{hyperref}
\hypersetup{colorlinks=true,linkcolor=blue!60!black,urlcolor=blue!60!black,citecolor=blue!60!black}
\usepackage{enumitem}
\setlist{nosep,leftmargin=1.6em}
\emergencystretch=1.5em  % texttt 长词与 CJK 混排时给断行弹性，消除边缘 overfull
\usepackage{fancyhdr}
\pagestyle{fancy}
\fancyhf{}
\fancyhead[L]{\small TeXlate 评测体系与指标演化}
\fancyhead[R]{\small 2026-09-19}
\fancyfoot[C]{\small\thepage}
\renewcommand{\headrulewidth}{0.4pt}

% 图配色
\definecolor{cA}{HTML}{1f6fb2}   % 主蓝
\definecolor{cB}{HTML}{c0392b}   % 红
\definecolor{cC}{HTML}{8a6d00}   % 棕
\definecolor{cD}{HTML}{2e8b57}   % 绿
\definecolor{cE}{HTML}{7d3c98}   % 紫
\definecolor{cF}{HTML}{16a085}   % 青

\title{\bfseries TeXlate 评测体系与指标演化全景报告\\[4pt]
  \large 七臂评测套件 · corpus\_v3 全量综合测试 · 2026-09-19}
\author{texlate bench 工作组}
\date{2026 年 9 月 19 日}

\begin{document}
\maketitle

\begin{abstract}
\noindent 本文是 TeXlate 评测体系的全景报告：上半部分为不熟悉项目的读者提供背景（管线架构、为什么 arXiv LaTeX 翻译难、评测套件设计与评测现场布局），下半部分是数据——全部历史 commit 考古出的指标演化时间线、2026-09-19 在 corpus\_v3 全量 13{,}266 篇上的综合测试画像、失败归因图谱与已知缺口，以及十条 takeaway。全部数字可追溯：每个指标注明口径与样本框；图表全部由随档 \texttt{data/} 原始数据生成，不依赖本目录以外的图片路径；\texttt{refs/} 仅随档上一期复盘的两张 SVG，引用源文档的时点快照未随迁——对应档案位置见附录「随档清单与原始产物索引」对照表及 README.md。
\end{abstract}

\vspace{1em}
\tableofcontents
\newpage

\section{背景知识}

\subsection{项目目标}

TeXlate 是「幻觉翻译」(hjfy.top) 的开源复刻：输入一篇 arXiv 论文的 LaTeX 源包，输出一份重编译的中文 PDF，正文逐段替换为中文译文、数学公式/图表/参考文献结构原样保留，并提供双语对照阅读。链路上的每一段都由自研组件构成：

\begin{itemize}
  \item \textbf{获取层}：arXiv 源包批量下载（IA 批量档 / export / OAI-PMH 三通道），钉版入库；
  \item \textbf{解析层}：自研半解析器——「逐字符扫描 $\to$ 切片段 $\to$ 占位符保护 $\to$ 区间回填」的保真切分器（非完整 TeX 引擎），外加宏展开机（Gullet）把用户自定义宏就地展开；
  \item \textbf{翻译层}：段落级译段经 LLM 网关批量翻译（swe-2 系列模型），占位符协议保护公式与命令不被译文污染；
  \item \textbf{校验层}：L0/L1/L2 三层校验（译段不变量、tree-sitter 结构校验、LLM 评审）；
  \item \textbf{编译层}：ctex 注入 + 规范化 + xelatex/tectonic 双引擎编译；
  \item \textbf{修复层}：修复循环（fixloop）——143 条 YAML 规则驱动（装包、改源、路由、polyfill），把编译失败的论文迭代救回；
  \item \textbf{对齐层}：PDF 具名锚点同步，滚动阅读时中英逐段对齐。
\end{itemize}

\subsection{主要难点}

\textbf{arXiv 源码是脏的真实世界数据}。语料横跨 LaTeX2.09（1990s）到 2026 年：混用 plain \TeX{} 原语与 LaTeX2.09 旧式、非 UTF-8 编码（Big5/GBK/Latin-1）、缺文件的残缺包、tar/gz 双格式、古早宏包（revtex4、aastex、psfig）。基线实测：裸编译（不做任何修复）能出 PDF 的仅 65--72\%，纯净出 PDF 的只有 22--40\%——三分之二的论文源码直接编译不出 PDF，最大失败类是缺文件（占裸编译失败的 92.1\%）。

\textbf{宏展开是现成解析器的共同盲区}。95\% 的论文含自定义宏（中位 47 个/篇，最多 446），其中 12/39 篇手挑语料存在「宏体里藏 \verb|\begin/\end|」的结构性陷阱，49\% 的宏体藏 $>20$ 字母的可译文本。对 8 个第三方解析库（latex-utensils、unified-latex、tree-sitter-latex、TexSoup、plasTeX、pylatexenc、LaTeX.js、ieeA）的横评（原始报告存档于 \texttt{bench/archive-2026-09-20/results/00-grand-comparison.md}）表明：陷阱断言全军覆没——最好成绩 24/26 vs 自研半解析器 32/32；「能解析」不等于「解析对」，pylatexenc 报 90/90 成功却是静默截断零报错——静默失败比显式崩溃更危险，故主解析器必须自研、校验器必须独立。

\textbf{翻译约束是结构性的}。译段边界纪律要求公式、命令、环境标签逐字节保留，译文只能在散文区间替换；术语表一致性、学术语域、不可译段（参考文献/版权行）都要显式协议。编译侧中文需要 ctex 注入与 CJK 字体路由，修复侧要处理 EPS 图源、缺 sty/cls、文献链路断点等 $\sim$260 类已建档机制。

\subsection{管线架构}

图~\ref{fig:pipeline} 给出全链架构与各段的 benchmark 对应关系。

\begin{figure}[htbp]
\centering
\resizebox{\textwidth}{!}{%
\begin{tikzpicture}[
  font=\small,
  stage/.style={draw,rounded corners=2pt,minimum height=8mm,inner xsep=4pt,align=center,fill=cA!8},
  impl/.style={font=\scriptsize\color{black!60},align=center},
  bench/.style={draw=cC,dashed,rounded corners=1pt,font=\scriptsize\color{cC},inner sep=2pt,align=center,fill=cC!6},
  arr/.style={-{Stealth[length=2mm]},thick,black!70},
  node distance=5mm and 4mm
]
% 主流水线
\node[stage] (fetch) {arXiv\\源包};
\node[stage,right=of fetch] (parse) {合平 + 半解析\\mouth/gullet/segmenter};
\node[stage,right=of parse] (xlat) {段落级翻译\\编排 + 网关};
\node[stage,right=of xlat] (valid) {校验\\L0/L1/L2};
\node[stage,right=of valid] (splice) {回填重建\\+ ctex 注入};
\node[stage,right=of splice] (comp) {编译\\xelatex / tectonic};
\node[stage,right=of comp] (align) {锚点对齐\\+ 双语 PDF};
\draw[arr] (fetch) -- (parse);
\draw[arr] (parse) -- (xlat);
\draw[arr] (xlat) -- (valid);
\draw[arr] (valid) -- (splice);
\draw[arr] (splice) -- (comp);
\draw[arr] (comp) -- (align);
% 修复循环回环
\node[stage,below=6mm of comp,fill=cB!10,draw=cB] (fix) {修复循环\\143 条 YAML 规则};
\draw[arr,cB] (comp.south) -- (fix.north);
\draw[arr,cB] (fix.west) -| (splice.south);
% 评测臂标注
\node[bench,below=4mm of parse] (b1) {B1 parsebench\\一致/泄漏/合平};
\node[bench,below=4mm of xlat] (b4) {B4 xlatbench + qualbench\\硬契约/延迟/ESA 评分};
\node[bench,below=11mm of valid] (b6) {B6 validbench\\检出率/误报};
\node[bench,below=4mm of fix] (b3) {B3 compilebench\\纯净/出PDF率};
\node[bench,below=4mm of align] (b7) {B7 alignbench\\锚点保留};
\node[bench,above=4mm of xlat] (b5) {B5 e2ebench 全链漏斗};
\draw[cC,dashed,-{Stealth[length=1.5mm]}] (b5.south) -- (xlat.north);
\node[bench,above=4mm of parse] (b2) {B2 fixtures T01--T98};
\end{tikzpicture}}%
\caption{TeXlate 管线架构与评测臂对应关系。主流水线自左向右；修复循环是编译段的自动修复回环；B1--B7 为七个评测臂（虚线框），另有 fixtures 陷阱断言集做解析段单元级测试。}
\label{fig:pipeline}
\end{figure}

\subsection{术语与口径}
\label{sec:terms}

后文所有指标沿用此表口径。

\begin{itemize}
  \item \textbf{格（cell）}：语料库中一篇论文的存储单元，目录为 \texttt{\{id\}/\{meta.json, raw.*, extracted/\}}；评测语境下也指一个被评对象（论文 $\times$ 引擎组合）。
  \item \textbf{层（layer）}：corpus\_v3 的抽样分层（见 \S\ref{sec:corpus}），每层独立清单、独立抽样框。\textbf{EVAL\_ONLY 层}为评测专用，开发期采样一律经 \texttt{dev\_layers()} 排除。
  \item \textbf{臂（arm）}：同一被测对象的不同处理条件——编译臂分裸编译（base）与中文链（zh，产品链规范化+ctex 注入后）；端到端/翻译臂分模拟（mock，确定性假翻译、零成本可复现）与真实（real，真 LLM 网关）；批量驱动的 base 为归因对照臂。编译条件记号：pipe-xel = 产品链后 xelatex 编译，pipe-fix = 修复循环后再编。
  \item \textbf{译段（chunk）与片段（piece）}：piece 为解析器输出的最小片段（散文/公式/命令等类型）；chunk 为翻译最小单位（段落级，由 piece 聚合而来），本批语料译段均长 211 字符。
  \item \textbf{合平（flatten）}：多文件论文（\verb|\input/\include| 树）合并为单一 .tex 的预处理；孤儿 tex 为不被主文档引用的附属文件（随附 preamble/poster 件），不计入合平覆盖。
  \item \textbf{回填（splice）与回填残留}：把译文按占位符区间替换回源文骨架的重建步骤；回填残留 = 重建后仍残留的占位符数，正常应为 0。
  \item \textbf{纯净（clean）/ 带瑕（pdf\~{}）/ 失败（FAIL）}：编译终态三级——纯净 = 无告警出 PDF；带瑕 = 带可接受瑕疵出 PDF（落盘记录记作 partial）；失败 = 无 PDF。\textbf{出 PDF 率 = 纯净 + 带瑕}。
  \item \textbf{联合口径（union）}：计分卡的出 PDF 统计是「阶段并集」——每格取编译与修复循环两段的较优终态；与之相对的终末口径只看最末段状态。base 臂是归因基线，不计入联合。
  \item \textbf{逐字节一致率（identity）与泄漏（leak）}：解析质量两轴——一致率为重建后逐字节相符的文件占比；泄漏为受保护内容（公式/命令）误入可译译段的占比。
  \item \textbf{死占位符（dead ph）与孤儿译段（orphan）}：解析侧两类零容缺——占位符未被任何译段消费称死占位符；译文译段找不到回填位置称孤儿。
  \item \textbf{批量驱动（stagerun）与计分卡（scorecard）}：stagerun 为五段批量驱动（获取$\to$解析$\to$翻译$\to$编译$\to$修复，逐格追加式落盘 records）；scorecard 为 \texttt{gate\_scorecard} 对记录的终态统计，M2 出口门出 PDF 率 $\geq 90\%$、纯净率 $\geq 90\%$。
  \item \textbf{硬契约（hard\_ok）与契约族}：翻译臂的硬契约判定——占位符全还原（\texttt{ph\_missing}、\texttt{ph\_invented}）、控制序列保留（\texttt{cs\_dropped}）、结构校验通过（\texttt{validator\_ok}）。顺序类分硬换序（\texttt{ord\_hard}，破契约）与软换序（\texttt{ord\_soft}，合法中文换序告警，不计败）。
  \item \textbf{ESA 评审与争议判例（contested）}：质量臂的 LLM 评审协议——评审模型按 ESA 错误分类打 0--100 分；初裁与复核分歧的判例记为争议，由第二评审仲裁。
  \item \textbf{具名锚点（named-dest）}：PDF 具名锚点（hyperref \verb|\label| 位点）；对齐臂测双版锚点保留率实现逐段对齐。
  \item \textbf{修复循环（fixloop）/ 机制（mechanism）/ 车道}：fixloop 为 YAML 规则驱动的编译修复循环；mechanism 为机制台账（\texttt{mechanisms.jsonl}，存档于 \texttt{bench/corpus/}）登记的已知缺陷模式（dollar 族泄漏、eps\_route 等），是「benchmark 发现 $\to$ 归因 $\to$ 规则沉淀」反馈环的载体；车道/波次为多代理并行修一批同族缺陷的工作单元；\textbf{自改进环（selfimp）}为 09-18 起常驻的多代理改进环（车队台账快照存档于 \texttt{tmp/old-docs-2026-09-20/docs/research/metrics-2026-09-19/refs/}）。
  \item \textbf{v2 切换（cutover）}：09-16 解析主路径从 v1 逐字符扫描器切到 v2 展开机 + 切分器的事件，切换硬门为回填残留占位符清零。
  \item \textbf{批量档（bg）}：网关批量档——低优先级队列，入闸排队约 120 秒，评测批跑默认档位，延迟读数含排队时间。
\end{itemize}
   % 背景知识 + 管线架构
\section{评测体系设计}

\subsection{七臂总览}

评测套件按「每臂对准管线一段 + 一个组合层」设计：每条臂独立成 harness，产出机读落盘记录供计分卡与归因消费。评测器规格冻结件为 \texttt{docs/spec/benchmark.md}，共同底材为 corpus\_v3。表~\ref{tab:arms} 为总览。

\begin{table}[htbp]
\centering\small
\caption{七臂评测总览（B1--B7 + 辅助件）}
\label{tab:arms}
\setlength{\tabcolsep}{4pt}
\begin{tabular}{@{}lllll@{}}
\toprule
\# & 评测臂 & 测哪段 & 底材 & 核心指标 \\
\midrule
B1 & parsebench & 解析段 & corpus\_v3 + corpus39 & 解析成功/一致率/泄漏/死件/合平 \\
B2 & fixtures & 解析段（单元级） & 手构陷阱 \texttt{.tex} 集（10 件） & 陷阱断言通过率（T01--T98） \\
B3 & compilebench & 编译段 + 修复循环 & corpus\_v3 原始包 & 纯净/带瑕/失败、救回率、规则谱 \\
B4 & xlatbench & 翻译段 & corpus\_v3 译段抽样 & 硬契约率 / 延迟 / token 经济 \\
B4b & qualbench & 翻译质量 & e2e 产物段对 & LLM 评审 ESA 评分 / 缺陷分类 \\
B5 & e2ebench & 全链（组合） & corpus\_v3 子集 & 逐段漏斗 + 终态分布 + 引入回归 \\
B6 & validbench & 校验段 & 合成破坏案例 & 检出率 / 误报 / 延迟 \\
B7 & alignbench & 阅读体验 & 中英编译产物对 & 具名锚点保留率 / 链权 \\
\midrule
\multicolumn{5}{@{}l}{辅助件：stagerun（五段批量流水线）、gullet\_bench（展开机实测）、}\\
\multicolumn{5}{@{}l}{gate\_scorecard（M2 出口门）、translators\_bench（模拟臂工厂）、nightwatch}\\
\bottomrule
\end{tabular}
\end{table}

依赖序：corpus\_v3 $\to$ B1 $\to$ B4 $\to$ \{B3, B5\} $\to$ B7；B2/B6 以合成输入独立。

\subsection{评测现场：\texttt{bench/} 的构成}

所有评测资产集中在仓库 \texttt{bench/} 树，入口文档为 \texttt{bench/TIERS.md}（验证分层契约）与 \texttt{docs/spec/benchmark.md} §3（逐库横评协议——原 \texttt{bench/PROTOCOL.md} 已于 2026-09-20 退役）。后者定义「一个问题该在哪一级被抓住」：L0 单测/断言样例（秒级，CI 硬门）$\to$ L1 语料机制覆盖（分钟级）$\to$ L2 stagerun 子集回归（分钟到小时，机制台账定向重放）$\to$ L3 全层全臂过夜集成（里程碑门）；规则是「能低不高」——低层能钉住的问题不许只在批量上撞见，批量发现的真坑要沉淀回 L0 层断言。本报告 B1--B7 属 L1--L3 层读数。布局与职能：

\begin{table}[htbp]
\centering\small
\caption{\texttt{bench/} 现场布局（不熟悉仓库的读者对照此表读结果目录名）}
\label{tab:benchlayout}
\begin{tabular}{@{}l>{\raggedright\arraybackslash}p{10.8cm}@{}}
\toprule
路径 & 内容 \\
\midrule
\texttt{bench/py/} & 37 个评测脚本，分四簇——\textbf{批量驱动}：\texttt{stagerun.py} + \texttt{stage\_\{ingest,parse,xlat,compile,fixloop\}}（五段流水线，落盘记录逐格追加式写入）；\textbf{七臂评测驱动}：\texttt{parsebench}、\texttt{compilebench\_v3}、\texttt{xlatbench}、\texttt{qualbench}、\texttt{e2e\_mock\_bench}、\texttt{e2e\_real\_bench}、\texttt{fixloop\_bench}、\texttt{alignbench}、\texttt{validbench}、\texttt{gullet\_bench}、\texttt{fixture\_assert}；\textbf{门禁监控}：\texttt{gate\_scorecard}（M2 出口门）、\texttt{status\_panel}、\texttt{nightwatch}、\texttt{rundiff}、\texttt{triage}、\texttt{wave}、\texttt{stage\_timing}；\textbf{支撑件}：\texttt{benchlib} 公共库、\texttt{translators\_bench}（模拟译文工厂）、\texttt{corpus/} 语料构建管线、\texttt{iclr\_*} 对照语料渠、\texttt{report/} 一次性审计、\texttt{scratch/} 一次性探针 \\
\texttt{bench/ts/} & JS 侧三库横评驱动（latex-utensils/unified-latex/tree-sitter-latex），D0 选型证据产地 \\
\texttt{bench/corpus*} & 语料族：\texttt{corpus} 39 篇手挑陷阱、\texttt{corpus\_v2} 139 篇分层、\texttt{corpus\_v3} 13{,}266 篇八层钉版（\S\ref{sec:corpus}）、\texttt{corpus\_daily} RSS 日更渠（$\sim$1{,}200 篇/日）、\texttt{corpus\_iclr*} ICLR 对照渠 \\
\texttt{bench/fixtures/} & 10 件手构陷阱 \texttt{.tex}（\texttt{@Tnn}/\texttt{@Wnn}/\texttt{@Xn} 断言标记：T 为结构陷阱、W 为野外发现机制、X 为翻译层；逐字节即语义，不做任何格式化） \\
\texttt{bench/results/} & 评测产物区：每次运行一个日期目录（377+ 轮累积），含 records/summary/jsonl/PDF 中间件；大体量不入档，本报告读数全部由其抽取进 \texttt{data/} \\
\texttt{bench/work\_*/} & 编译/fixloop 逐格工作区（解压源树、aux、PDF；重产物 gitignored） \\
\bottomrule
\end{tabular}
\end{table}

\subsection{语料分层构造}
\label{sec:corpus}

corpus\_v3 是钉版批量语料：每格记录 \texttt{(channel, item, member, blob\_sha256)} 四元组钉住字节级来源（\texttt{resolved\_version=null}），落地为 \texttt{\{id\}/} 目录（\texttt{meta.json}、\texttt{raw.*}、\texttt{extracted/}）。四条渠道各有明确分工：\texttt{ia} = IA \texttt{arxiv-bulk} 月度块（a--d 时代带主力，覆盖至 2020-10）；\texttt{tiger} = HF \texttt{TIGER-Lab/arxiv-latex-5T}（e 带主力，成员截止 2412 月）；\texttt{arxiv\_eprint} = 单发源包取源（走产品 \texttt{acquire\_source}，补 TIGER 截止后 2501+ 盲区）；\texttt{hf\_scholarweave} = scholarweave 脱水通道（\texttt{figures\_stripped} 有损源——解析/翻译段可用，编译段因缺图源排除）。全库 13{,}266 格、原始包 18.8GB（含 \texttt{extracted/} 解压树全库约 46GB）。八层构成见图~\ref{fig:corpus}、分层定位见表~\ref{tab:layers}：核心/补强/扩展/热层为早期层，2026-09-19 评测/开发分轨扩层后收至 13{,}266 篇。

\begin{figure}[htbp]
\centering
\begin{tikzpicture}
\begin{axis}[
  xbar stacked,
  width=13cm, height=4.6cm,
  xmin=0, xmax=14000,
  scaled x ticks=false,
  xlabel={千篇}, ytick=data,
  xtick={0,2000,4000,6000,8000,10000,12000,14000},
  xticklabels={0,2,4,6,8,10,12,14},
  yticklabels={corpus\_v3 八层},
  legend style={at={(0.5,-0.28)},anchor=north,legend columns=4,font=\scriptsize},
  bar width=8mm,
  x label style={font=\small},
]
\addplot+[fill=cA!75,draw=cA] coordinates {(1000,0)};
\addplot+[fill=cB!75,draw=cB] coordinates {(200,0)};
\addplot+[fill=cD!75,draw=cD] coordinates {(3866,0)};
\addplot+[fill=cE!75,draw=cE] coordinates {(166,0)};
\addplot+[fill=cF!75,draw=cF] coordinates {(2000,0)};
\addplot+[fill=cC!75,draw=cC] coordinates {(1514,0)};
\addplot+[fill=cB!40,draw=cB] coordinates {(1500,0)};
\addplot+[fill=black!35,draw=black!60] coordinates {(3020,0)};
\legend{core 1000, booster 200, expand 3866, hot 166, dev\_vol 2000, dev\_recent 1514, dev\_failmine 1500, holdout 3020}
\end{axis}
\end{tikzpicture}
\caption{corpus\_v3 八层构成（13{,}266 篇，原始包 18.8GB / 含解压树约 46GB）。holdout 层为 EVAL\_ONLY 治理——\texttt{dev\_layers()} 实测排除，仅评测可读；dev\_* 三层与 expand 为开发期增量；core 为月度簇分层随机（30 簇，old 200/new 800）。}
\label{fig:corpus}
\end{figure}

\begin{table}[htbp]
\centering\small
\caption{corpus\_v3 八层定位与抽样框（渠道缩写：ia=IA 月块 / tiger=HF arxiv-latex-5T / eprint=单发取源 / sw=scholarweave 脱水）}
\label{tab:layers}
\begin{tabular}{@{}ll>{\raggedright\arraybackslash}p{3.4cm}>{\raggedright\arraybackslash}p{6.4cm}@{}}
\toprule
层 & 篇数 & 定位 & 渠道构成与抽样框 \\
\midrule
core & 1{,}000 & 核心随机层，评测主干 & ia 800 + tiger 200；30 个 IA 月簇分层随机（old 200 / new 800） \\
booster & 200 & 机制策展层 & ia 162 + tiger 38；野外发现机制提名制（107/109 W exemplars，\texttt{nominations/} 留审计轨迹） \\
expand & 3{,}866 & 体量扩展层 & ia 3{,}103 + tiger 697 + eprint 66；bulk 池 measure-then-sample \\
hot & 166 & 近期高引层 & arxiv\_eprint 166；OpenAlex 高引近期，专打 2501+ 盲区 \\
dev\_vol & 2{,}000 & 开发体量层 & ia 1{,}627 + tiger 373；3{,}671 tex $\cdot$ 2.76GB \\
dev\_recent & 1{,}514 & 近期开发层 & sw 1{,}065 + eprint 449；TIGER 截止后盲区双通道 \\
dev\_failmine & 1{,}500 & 机制挖掘层 & ia 1{,}500；FLAG\_RX 机制旗标定向挖 IA 94--15 年档陷阱 \\
holdout & 3{,}020 & 留出评测层 & ia 2{,}159 + tiger 540 + eprint 321；EVAL\_ONLY，同构造规则独立抽样 \\
\bottomrule
\end{tabular}
\end{table}

抽样框是「时代带 + 机制旗标带」的二维结构：五个时代带 a\_pre2007 / b\_2007\_11 / c\_2012\_16 / d\_2017\_20 / e\_2021\_25 各 $\sim$1{,}700--2{,}050 篇均匀打底（时间跨度 1994--2026），上面叠机制旗标带定向覆盖已知陷阱族（\texttt{flag:deadpkg} 600、\texttt{flag:docstyle209} 300、\texttt{flag:epsfig}、\texttt{flag:babel}、\texttt{flag:pdftex}、\texttt{flag:pstricks} 等若干）与 \texttt{sw|yymm} 月度带 1{,}065。学科面随 arXiv 主力盘分布（图~\ref{fig:corpusdist} 左）：hep-phys 2{,}961、math 2{,}418、cs 2{,}393、astro-ph 1{,}826、cond-mat 1{,}729 为五大盘。

物理形态与工程结构：tar 包 11{,}097、gz 单文件 2{,}144、single\_gz 12、stub 13（B07 有意构造的残缺件）；单篇文件数中位 7 / p90 28 / 最大 3{,}047（一篇论文三千文件），\texttt{.tex} 均 2.2 个，多文件工程 2{,}115 篇（16\%）——合平与多文档路由不是边角料而是主场景之一。许可面：arxiv-nonexclusive 6{,}786、missing 3{,}335、CC 系合计 1{,}187。全库 70 个聚簇键（\texttt{cluster\_id}，月簇/渠道混合）支撑簇稳健自举；4{,}472 格带 B01--B07/W/T 机制标签——即「评测发现 $\to$ 归因 $\to$ 策展回填」反馈环的载体。

\begin{figure}[htbp]
\centering
\begin{tikzpicture}
\begin{groupplot}[
  group style={group size=2 by 1,horizontal sep=18mm},
]
\nextgroupplot[
  xbar, bar width=3.4mm,
  width=7.6cm, height=5.2cm,
  xmin=0, xmax=3450,
  symbolic y coords={nucl,quant-ph,eess-stat,cond-mat,astro-ph,cs,math,hep-phys},
  ytick=data, y tick label style={font=\scriptsize},
  xlabel={篇数}, x label style={font=\small},
  nodes near coords, nodes near coords style={font=\tiny},
]
\addplot+[fill=cA!70,draw=cA] coordinates {(282,nucl) (552,quant-ph) (583,eess-stat) (1729,cond-mat) (1826,astro-ph) (2393,cs) (2418,math) (2961,hep-phys)};
\nextgroupplot[
  ybar, bar width=4.4mm,
  width=7.2cm, height=5.2cm,
  ymin=0, enlarge y limits={upper,value=0.22},
  symbolic x coords={a$<$07,b 07--11,c 12--16,d 17--20,e 21--25,sw,flag},
  xtick=data, x tick label style={font=\tiny,rotate=28},
  ylabel={篇数}, y label style={font=\small},
  nodes near coords, nodes near coords style={font=\tiny},
]
\addplot+[fill=cD!70,draw=cD] coordinates {(a$<$07,2039) (b 07--11,2047) (c 12--16,1936) (d 17--20,1726) (e 21--25,1817) (sw,1065) (flag,2636)};
\end{groupplot}
\end{tikzpicture}
\caption{corpus\_v3 构成两视图。左：学科组分布 top8（其余 470 篇未绘）；右：stratum 带分布——a--e 五时代带共 9{,}565 篇均匀打底，sw 月带 1{,}065，flag 机制带/dev 补/holdout recent 等合计 2{,}636。}
\label{fig:corpusdist}
\end{figure}

治理规则：holdout 层（3{,}020 篇）登记 \texttt{EVAL\_ONLY\_LAYERS}，所有开发期采样经 \texttt{dev\_layers()} 过滤；语料扩展只加层不删层；格数据不入库（gitignored），入库物为清单八件与选样报告（\texttt{selection\_report.md}）；机制台账 \texttt{mechanisms.jsonl}、B04/B06 宇宙×语料覆盖簿记 \texttt{eval\_coverage.json}、补强层提名审计轨迹 \texttt{nominations/} 均为 \texttt{bench/corpus/} 下正档——台账是「评测发现 $\to$ 归因 $\to$ 规则沉淀」反馈环的载体。

\subsection{逐臂详解}
\label{sec:arms}

\subsubsection{B1 parsebench —— 解析正确性}

\textbf{测什么}：\texttt{texlate.latex} 半解析器在真实语料上的正确性——能不能零崩溃跑完、能不能逐字节还原、可译译段里有没有漏进受保护内容。

\textbf{怎么测}：逐文件 30 秒超时，先定位主文件，再做多文件合平，随后调用解析器，最后跑重建三连测——逐字节一致（strict identity）、空译文假翻译（fake-translate，注入死占位符/孤儿探针）、占位译文（泄漏探针）。逐文件结果写入 \texttt{files.jsonl}，逐篇聚合成 \texttt{papers.json}，同时产出路由标签（reject / xelatex / minted / non-utf8 / no-hyperref——这组标签同时充当 B3 的静态路由金标准）以及合平覆盖率、多文档/无根标记。统计口径为核心层加权池化率与宏平均双报，置信区间用 Wilson 加月簇稳健自举；另有 oracle 对拍探针（plasTeX 规范化文本流对译段的字符级召回），仅作开发期召回探针、不进门槛。所有异常逐条人工归因后写回机制台账。

\textbf{当前读数}：28{,}904 文件解析成功率 100.0\%、逐字节一致率 99.99\%、泄漏率 0.004\%（图~\ref{fig:parse}、表~\ref{tab:snapshot}）。

\subsubsection{B2 fixtures —— 机制断言集}

\textbf{测什么}：parsebench 测分布，断言集测机制——「这个具体的坑处理了吗」逐条断言。

\textbf{底材}：10 件手构 \texttt{.tex}（tricky 主集 + 209 旧格式件 + multi 多文件工程件 + w/w73/wenc 野外机制件 + xlat-traps 翻译层件），\texttt{@Tnn}（结构陷阱）/\texttt{@Wnn}（野外发现机制）/\texttt{@Xn}（翻译层）三族标记；\textbf{逐字节即语义}，永不格式化。\textbf{生长机制}：parsebench 归因与台账中 \texttt{found-in-wild} 条目达到 covered 状态后提取固化为断言样例——语料里每个真坑沉淀为一条永久断言，断言面只增不减。

\textbf{当前读数}：98 条断言 pytest 全绿（L0 层 CI 硬门，断言函数与 \texttt{fixture\_assert.py} 共享）；BUG1 占位符尾回归断言实测为 0。

\subsubsection{B3 compilebench —— 编译与修复循环救回}

\textbf{测什么}：规范化 + ctex 注入 + 引擎 + 修复循环的真实救回能力——hjfy 约五千篇人肉沉淀护城河的对应物。

\textbf{怎么测}：论文 $\times$ 条件 $\times$ 引擎三维网格——条件取原文裸编（base）/ 中文注入（zh）/ 模拟译文（mock）三种，引擎取 xelatex 与 tectonic 两种；编译前先做路由预检（\texttt{documentstyle}$\to$直接拒编、eps/pstricks$\to$xelatex、minted-frozencache$\to$tectonic）。每格走「规范化 $\to$ 中文注入 $\to$ 沙箱编译（至多 2 轮、240 秒限时）$\to$ 修复循环 $\to$ 纯净判定三件套」；逐格判定（verdict）连同每轮 \texttt{\{cat,pay,rule,result\}} 动作落盘，直接灌回机制台账。失败归因按首错签名聚类（\texttt{verdict\_sig} 口径：当构成中计数最高的错误类超过首错类时改挂该众数类）。指标：纯净率（分条件/引擎/时代带）、失败$\to$出 PDF 救回率、规则触发/救回数、卡死率、修复轮数分布。

\textbf{门槛}：200 篇中文链条件编译成功率 $\geq$90\%（hjfy 的 95\% 为渐近线）、拒编判定正确率 100\%（与路由标签对拍）、无回归（原本纯净的格不被规则改脏）。

\textbf{当前读数}：500 格联合口径出 PDF 90.4\%（图~\ref{fig:compile}）；中文臂 xelatex 纯净率 +16.1 个百分点——规范化顺带修复源缺陷，产品链成本转正收益。

\subsubsection{B4 xlatbench / qualbench —— 翻译契约与质量}

\textbf{B4a 硬契约层}：译段 $\times$ 模型 $\times$ 重复次数网格，译段按类型分层抽样（正文段、图题、条目、脚注，外加压力样例——\verb|\bibitem| 前缀、\verb|\href| 等陷阱位）；每个响应过 L0 校验器，指标按序读——硬契约通过率（丢/造占位符、丢脆弱命令、校验错误）$\to$ 换序软信号（合法中文换序不降权）$\to$ 延迟 $\to$ 推理量 $\to$ token 经济。用途是模型选型与白名单刷新（免费集推广期到期即重跑）以及提示词措辞回归（每次提升 \texttt{prompt\_version} 必跑）。

\textbf{B4b qualbench 质量层}：LLM 评审（esa2 协议）逐段对打 0--100 分，并按错误分类 $\times$ 严重度三级打标（术语/流畅/准确/约定/漏译五族九类）；初裁 swe-2-max，分歧判例记为争议，由 swe-2-high 复核仲裁；产出均分/中位分、分数段分布、按译段类型与按论文的细分。门槛 M1：100 篇真实翻译端到端 $\geq$85\%。

\textbf{当前读数}：硬契约通过率 93.7\%（254/271）、ESA 均分 94.0 / 争议率 5.5\%（图~\ref{fig:qual}）。

\subsubsection{B5 e2ebench —— 全链组合}

\textbf{测什么}：单段绿不等于组合绿——逐环节成功率漏斗加终态分布。四模式：\textbf{Mode A} 位置忠实的模拟译文（用占位译文走全链验证机械链路：能否出 PDF、逐字节一致、零残留占位符；已知盲区是模拟翻译只认 ASCII 散文段，非 ASCII 散文会原样回显、过估忠实度——真实翻译臂不受影响）；\textbf{Mode B} 幻觉模拟（注入占位符丢失/幻觉破坏，验证校验链兜底——132 处破坏编译前 132/132 捕获，「出 PDF $\neq$ 成功」的实证来源）；\textbf{Mode C} 位置扰动模拟（随机移动约 10\% 占位符，量化回填鲁棒性）；\textbf{Mode D} 真实翻译接入（产品 SLA 观测点）。另有 \textbf{pipe-fix 救回臂}：把 pipe-xel 的产物树送进修复循环复判，联合口径取两臂较优者；\texttt{onfail} 语义经冒烟校准为「失败 + misschar/error 级带瑕」才救（带瑕已出 PDF，引擎重跑只会丢 PDF 救不了告警），注入拒编不救。

\textbf{当前读数}：模拟臂管线 xel 33/39 纯净 vs 裸编 19/39（管线净正 +14 格）；真实臂 60 篇联合口径出 PDF 96.7\%、回填残留 0、0 管线引入回归。

\subsubsection{B6 validbench —— 校验器自检}

\textbf{测什么}：L0/L1 校验器对 LLM 破坏的检出能力——校验器本身必须被评测（实现期真抓到过校验器自身 bug）。

\textbf{怎么测}：变异器对干净译段施加 10 类破坏（丢右括号、丢美元符、\verb|\end| 改名或删除、丢/拼错/凭空多占位符、\verb|\[| 不配对、幻觉宏、删引用键），在占位符层/原文层两层生成 \texttt{cases.jsonl}；另加对抗性手工探针（丢 \verb|\emph| 组、占位符合法换序、\verb|\cite {k}| 带空格、漏引用键、全角占位符等——验证合法改写不误报）。指标：检出率（按破坏类分列命中规则）、干净对误报率、仅告警率、每对延迟、编辑距离 $\leq$2 修复建议率；L1 tree-sitter 层同口径。B4 产出的真实失败案例反向喂回 B6，沉淀为新破坏类。

\textbf{门槛}：10 类破坏 100\% 检出、干净对 0 误报、L0 $\leq$1ms/对。

\textbf{当前读数}：corpus\_v2 1{,}503 对全绿（摊薄 0.374ms/对）——本批未重跑，列缺口清单第 7 项。

\subsubsection{B7 alignbench —— 锚点保留}

\textbf{测什么}：中文重编译后 PDF 具名锚点（named-destination）保留率——双语滚动同步的质量上限。

\textbf{怎么测}：用 pypdf 提取中英两侧具名锚点，同名锚点配对，计算保留率与最大权值单调链权重（节标题锚权 12、图锚 10、公式锚 4、引用锚 2，\texttt{page.*} 页锚权 0）。未用 hyperref 的工程（语料中约 31\%）走退化路径、只断言不崩溃；\textbf{保留率 $<$95\% 本身可当中文编译完整性探针}（ctex 共享计数器改 theorem 锚名的问题即由此捕获）。

\textbf{门槛}：有 hyperref 对保留率 $\geq$95\%；无 hyperref 对退化路径不崩。

\textbf{当前读数}：812 对门全过但覆盖薄——765 对无具名锚点，有效读数仅 42 个 hyperref 对（保留率中位 1.0），口径待复核（缺口清单第 3 项）。

\subsubsection{辅助件}

stagerun 五段批量流水线（逐格追加式落盘、按 \texttt{(id,arm,upstream)} 断点续跑、网关信号量限流），配 \texttt{gate\_scorecard} 做 M2 出口门（联合口径 = 编译 $\cup$ 修复两段逐格取优）；\texttt{translators\_bench} 为模拟臂译文工厂（忠实模拟 / 注入破坏 b、c / 位置扰动）；\texttt{gullet\_bench} 为展开机实测（200 篇全展开中位 39.6ms/19 步）；另有 \texttt{triage}（落盘记录 $\to$ 问题工单聚类）、\texttt{status\_panel}、\texttt{nightwatch}、\texttt{rundiff}、\texttt{wave}、\texttt{stage\_timing}、\texttt{preflight\_batch} 等监控与批前闸件。

\subsection{门槛体系}

各臂门槛（\texttt{docs/spec/benchmark.md} §8 门槛汇总）：parsebench——解析成功率 100\%、逐字节一致率 $\geq 99.5\%$、泄漏率 $\leq 0.15\%$、死占位符与孤儿译段 $=0$、合平覆盖率 $\geq 99\%$；M2 出口门（stagerun 计分卡）——联合口径出 PDF 率 $\geq 90\%$、纯净率 $\geq 90\%$。统计口径：核心层加权池化率 + 宏平均双报，置信区间用 Wilson 加月簇稳健自举。
        % 评测体系设计（七臂 + 语料 + 门槛）
\section{指标演化史}

本节数据来自三路考古：全部 1{,}648 条历史 commit 逐条扫描（里程碑识别）、评测结果流水 377 个日期目录的读数回填、车队台账与上一期复盘（口径锚点）。本文档自包含：原始数据入库于随档 \texttt{data/} 子目录（scorecard/parse/corpus/rules-tests/qual-xlat 五条时间线 CSV + 377 行评测活动流水含各轮 run 目录名 + 1{,}648 行 commit 全录 + milestones.md），引用源文档时点快照在随档 \texttt{refs/}（台账、复盘、评测器与语料规格、8 库横评报告、机制台账）；口径不一的读数分线绘制并注明样本框。

\subsection{开发节奏：六天 1{,}648 个 commit}

\begin{figure}[htbp]
\centering
\begin{tikzpicture}
\begin{axis}[
  width=12.5cm, height=5cm,
  ybar, bar width=9mm,
  ymin=0, enlarge y limits={upper,value=0.15},
  symbolic x coords={09-14,09-15,09-16,09-17,09-18,09-19},
  xtick=data, x tick label style={font=\small},
  ylabel={commits/日}, y label style={font=\small},
  nodes near coords, nodes near coords style={font=\scriptsize},
]
\addplot+[fill=cA!70,draw=cA] coordinates {(09-14,9) (09-15,123) (09-16,306) (09-17,624) (09-18,265) (09-19,321)};
\end{axis}
\end{tikzpicture}
\caption{commit 日吞吐（累计 1{,}648）。09-17 峰值 624 对应「波次战」（离线宏包 shim/\texttt{illegal\_unit}/段修多车道并行收割）；全程为多代理车队协作模式（监督台账记名册峰值约 23 车道），commit 由主会话统一落盘。}
\label{fig:commits}
\end{figure}

项目节奏呈「日主题」结构（图~\ref{fig:commits}、表~\ref{tab:days}）：D0 立项与八库横评定案、D1 产品骨架（七臂同日落地）、D2 v2 切换 + 修复循环接线（最大单日指标跃升）、D3 波次战收割残面、D4 自改进常驻环、D5 语料收官 + 全量综合测试。

\begin{table}[htbp]
\centering\small
\caption{六日大事记（commit 考古，精选）}
\label{tab:days}
\begin{tabular}{@{}ll>{\raggedright\arraybackslash}p{9.2cm}@{}}
\toprule
日 & 主题 & 标志事件 \\
\midrule
D0 09-14 & 立项定案 & 仓库初始化；\textbf{8 库横评}（39 篇 + 30 断言样例）——miniscanner 陷阱 32/32、逐字节一致 259/259 定案自研路线；技术栈选型冻结 \\
D1 09-15 & 产品骨架 & miniscanner 重写为 \texttt{texlate/latex} 九文件；Mouth+Gullet 展开机；编译/翻译/服务端/阅读器全套；\textbf{七臂同日落地}；修复循环 25 条规则移植；语料 $\to$1{,}339 \\
D2 09-16 & 切换+接线 & \textbf{v2 切换}（回填残留 1524$\to$0）；\textbf{修复循环接线}（出 PDF $+18.9$pt）；expand+hot 层语料 $\to$5{,}410；loop1 $n$=5124 计分卡 97.89\%；qualbench 诞生 \\
D3 09-17 & 波次战 & 离线宏包 23 件翻 492 格；\texttt{illegal\_unit} 族 108 闭环；ReDoS 修复 pytest 48min$\to$3:46；翻译层 P0 封死错配通道；\textbf{realn200 98.5\% 真臂验收} \\
D4 09-18 & 自改进环 & 常驻改进环 23 车道发车；loop2 97.26\%/86.42\%；\textbf{dollar 族泄漏清零}；qualbase 92.9 建档；corpus\_daily 日更渠启动 \\
D5 09-19 & 收官+全量 & \textbf{corpus\_v3 八层 13{,}266 篇}（开发/评测分治）；完整性审计全绿；\textbf{12 臂综合测试同日落地}（本报告主体）；loop3 在飞 \\
\bottomrule
\end{tabular}
\end{table}

\subsection{编译健康度}

\begin{figure}[htbp]
\centering
\begin{tikzpicture}
\begin{axis}[
  width=14cm, height=8cm,
  date coordinates in=x,
  xtick={2026-09-15,2026-09-16,2026-09-17,2026-09-18,2026-09-19},
  x tick label style={font=\scriptsize},
  ymin=0, ymax=108,
  ylabel={\%}, y label style={font=\small},
  legend style={at={(0.97,0.03)},anchor=south east,font=\scriptsize,legend columns=2},
  grid=major, grid style={black!15},
]
% scorecard union pdf（stagerun records 重评）
\addplot+[color=cA,mark=*,thick] coordinates {
  (2026-09-16,89.17) (2026-09-17,95.99) (2026-09-17,97.35)
  (2026-09-17,97.89) (2026-09-18,97.26) (2026-09-19,98.75)
};
% scorecard clean
\addplot+[color=cB,mark=*,thick] coordinates {
  (2026-09-16,57.20) (2026-09-17,68.91) (2026-09-17,82.76)
  (2026-09-17,84.99) (2026-09-18,86.42) (2026-09-19,88.75)
};
% real 臂
\addplot+[color=cE,mark=triangle*,mark size=3pt] coordinates {
  (2026-09-17,98.5) (2026-09-19,96.7)
};
\addplot+[color=cE,mark=triangle,mark size=3pt,dashed] coordinates {
  (2026-09-17,86.5) (2026-09-19,86.7)
};
% 裸编译基线
\addplot+[color=black!55,mark=square,only marks] coordinates {
  (2026-09-15,71.7) (2026-09-16,70.6) (2026-09-16,71.3) (2026-09-19,90.2)
};
\addplot+[color=cD,dashed] coordinates {(2026-09-15,90) (2026-09-19,90)};
\legend{联合口径出 PDF（计分卡）, 纯净（计分卡）, 真实臂联合出 PDF, 真实臂纯净, 裸编译联合出 PDF, M2 门 90\%}
\end{axis}
\end{tikzpicture}
\caption{编译健康度时间线。计分卡线为 stagerun 落盘记录在代码演进下的重评（loop1 $n$=5124 $\to$ loop2 $n$=5219 $\to$ v3all 样本 $n$=80）；真实臂为真 LLM 翻译全链（realn200 $\to$ e2e-v3b $n$=60 终态）；方块为裸编译基线——注意\textbf{基线自身四天涨约 19 个百分点}（71.7$\to$90.2）：离线宏包资产 + tlmgr usertree 积累使工具链整体成熟，「裸引擎撑死七成」的 09-15 结论已不适用今日地板。}
\label{fig:health}
\end{figure}

三条叙事线（图~\ref{fig:health}）：\textbf{修复循环是单一最大杠杆}——09-16 接产品链当日联合口径出 PDF $+18.9$ 个百分点（70.6$\to$89.5，同 180 篇复跑）；\textbf{波次战把残面打成个位数}——离线宏包波翻 492 格失败面、\texttt{illegal\_unit} 族 108/108 闭环、计分卡一波 89.17$\to$97.89；\textbf{真实臂与模拟臂打平}——realn200 联合出 PDF 98.5\% 证明评测面非虚高，今日 e2e-v3b 终态 96.7\%（7 个管线失败格全被修复循环救回、2 格超大工程跳过）。未达标项只剩 M2 纯净率 $\geq90\%$（今 88.75\%，差 1.25 个百分点）。

\subsection{解析侧：一致率守住，泄漏四天内归零再现身}

\begin{figure}[htbp]
\centering
\begin{tikzpicture}
\begin{groupplot}[
  group style={group size=2 by 1,horizontal sep=16mm},
  width=7.4cm, height=5cm,
]
\nextgroupplot[
  ybar, bar width=8mm,
  ymin=0, enlarge y limits={upper,value=0.2},
  symbolic x coords={09-15 核,09-17 核,09-18 核,09-19 全},
  xtick=data, x tick label style={font=\tiny,rotate=18},
  ylabel={泄漏命中数}, y label style={font=\small},
  nodes near coords, nodes near coords style={font=\scriptsize},
]
\addplot+[fill=cB!70,draw=cB] coordinates {(09-15 核,58) (09-17 核,57) (09-18 核,0) (09-19 全,76)};
\nextgroupplot[
  ybar, bar width=8mm,
  ymin=99.9, ymax=100.02,
  symbolic x coords={核心,补强,expand,全量},
  xtick=data, x tick label style={font=\tiny,rotate=18},
  ylabel={一致率 \%}, y label style={font=\small},
  ytick={99.9,99.95,100},
  nodes near coords, nodes near coords style={font=\scriptsize},
]
\addplot+[fill=cA!70,draw=cA] coordinates {(核心,100) (补强,100) (expand,100) (全量,99.99)};
\end{groupplot}
\end{tikzpicture}
\caption{解析侧演化。左：泄漏命中数——dollar 族 58/57 条残留于 09-18 c2-dollar 车道清零（核心层口径），09-19 全量 28{,}904 文件上重现 76 条（1{,}845{,}338 译段中占 0.004\%，新语料表面新案例，非回归）。右：逐字节一致率全测量期 $\geq$99.99\%（纵轴截断）。}
\label{fig:parse}
\end{figure}

解析侧四天稳态（图~\ref{fig:parse}）：v2 展开机+切分器切换当日回填残留占位符 1524$\to$0；泄漏残留清一色 dollar 族，09-18 单类收口归零——但\textbf{归零是在 1{,}955 文件核心层上}，今日 8 倍体量全量复测重现 76 条命中，说明该族机制在长尾语料上未穷举，属于「已识别族的覆盖率问题」而非新机制。

\subsection{翻译侧：硬契约与质量评分}

\begin{figure}[htbp]
\centering
\begin{tikzpicture}
\begin{groupplot}[
  group style={group size=2 by 1,horizontal sep=16mm},
  width=7.4cm, height=5cm,
]
\nextgroupplot[
  ybar, bar width=7mm,
  ymin=55, ymax=105,
  symbolic x coords={smoke 48,契约 80,全层 271},
  xtick=data, x tick label style={font=\tiny,rotate=18},
  ylabel={硬契约率 \%}, y label style={font=\small},
  nodes near coords, nodes near coords style={font=\scriptsize},
]
\addplot+[fill=cA!70,draw=cA] coordinates {(smoke 48,64.6) (契约 80,98.8) (全层 271,93.7)};
\nextgroupplot[
  ybar, bar width=7mm,
  ymin=88, ymax=96.6,
  symbolic x coords={qualbase,qualdrift,e2ereal},
  xtick=data, x tick label style={font=\tiny,rotate=18},
  ylabel={评审均分}, y label style={font=\small},
  nodes near coords, nodes near coords style={font=\scriptsize},
]
\addplot+[fill=cD!70,draw=cD] coordinates {(qualbase,92.9) (qualdrift,92.8) (e2ereal,94.0)};
\end{groupplot}
\end{tikzpicture}
\caption{翻译侧演化。左：swe-2-medium 硬契约率——09-15 冒烟期网关不稳（http\_err 17）压到 64.6\%，契约批 98.8\%，今日全层抽样 93.7\%（新校验口径 v\_err 27 条计入后更严）；右：ESA 评审均分 92.9$\to$94.0（swe-2-max 初裁，争议率 8.1\%$\to$5.5\% 下降）。}
\label{fig:xlatqual}
\end{figure}

\subsection{评测资产与代码规模}

\begin{figure}[htbp]
\centering
\begin{tikzpicture}
\begin{groupplot}[
  group style={group size=4 by 1,horizontal sep=8mm},
  width=4.1cm, height=4.8cm,
  enlarge y limits={upper,value=0.22},
]
\nextgroupplot[
  ybar, bar width=3.2mm, ymin=0,
  symbolic x coords={15,16,17,18,19},
  xtick=data, x tick label style={font=\tiny},
  title={pytest 用例}, title style={font=\tiny},
  nodes near coords, nodes near coords style={font=\tiny,rotate=90},
]
\addplot+[fill=cA!70,draw=cA] coordinates {(15,798) (16,1920) (17,4333) (18,5261) (19,6450)};
\nextgroupplot[
  ybar, bar width=3.2mm, ymin=0,
  symbolic x coords={15,16,17,18,19},
  xtick=data, x tick label style={font=\tiny},
  title={修复规则数}, title style={font=\tiny},
  nodes near coords, nodes near coords style={font=\tiny,rotate=90},
]
\addplot+[fill=cB!70,draw=cB] coordinates {(15,31) (16,50) (17,71) (18,106) (19,143)};
\nextgroupplot[
  ybar, bar width=3.2mm, ymin=0,
  symbolic x coords={15,16,17,18,19},
  xtick=data, x tick label style={font=\tiny},
  title={源码文件数}, title style={font=\tiny},
  nodes near coords, nodes near coords style={font=\tiny,rotate=90},
]
\addplot+[fill=cE!70,draw=cE] coordinates {(15,72) (16,89) (17,174) (18,355) (19,425)};
\nextgroupplot[
  ybar, bar width=3.2mm, ymin=0,
  symbolic x coords={c39,v2,+cor,+boo,+hot,+exp,+dv4},
  xtick=data, x tick label style={font=\tiny,rotate=45},
  title={语料篇数 /千篇（v3 分层累加）}, title style={font=\tiny},
  nodes near coords, nodes near coords style={font=\tiny,rotate=90},
]
\addplot+[fill=cD!70,draw=cD] coordinates {(c39,0.039) (v2,0.139) (+cor,1.0) (+boo,1.2) (+hot,1.366) (+exp,5.232) (+dv4,13.266)};
\end{groupplot}
\end{tikzpicture}
\caption{资产增长（git 历史抽样，同日口径一致）：pytest 用例 798$\to$6{,}450、修复规则 31$\to$143（16 分片）、源码 72$\to$425 文件、语料 corpus39/v2 $\to$ v3 八层 13{,}266（另有 corpus\_daily 日更渠约 1{,}200 篇/日在库外增长）。}
\label{fig:assets}
\end{figure}

\subsection{归因表：哪些改动搬动了哪些指标}

\begin{table}[htbp]
\centering\small
\caption{改动 $\to$ 指标位移归因（commit 级，摘自台账与各波报告）}
\label{tab:attrib}
\begin{tabular}{@{}ll>{\raggedright\arraybackslash}p{7.5cm}@{}}
\toprule
日期 & 改动（commit） & 指标位移 \\
\midrule
09-14 & 8 库横评 + miniscanner 探针 & 主解析器路线定案（陷阱 32/32，逐字节一致 259/259） \\
09-15 & corpus\_v3 core+booster + 双引擎基线 & 基线口径建立（联合出 PDF 71.7\%/纯净 39.4\%，180 篇） \\
09-15 & 修复循环 yaml 引擎落地（\texttt{179f7e5b}） & 25 条规则移植，修复臂诞生 \\
09-16 & 修复循环接产品链（\texttt{3d5de8f7}） & 联合出 PDF 70.6$\to$89.5\%（$+18.9$pt，同口径复跑） \\
09-16 & v2 展开机+切分器切换（\texttt{f4616838}） & 回填残留占位符 1524$\to$0（切换硬门） \\
09-16 & argspec.json 入包（\texttt{401e9bbf}） & 1820 条宏参规格零装载缺口闭合 \\
09-16 & perf-tail 五修 & 解析长尾 $-48\%$（51476$\to$26723ms），一致率零回退 \\
09-16 & stagerun loop1 全流水线（$n$=5124） & 计分卡 97.89\%/84.99\%；修复救回 84.0\% \\
09-17 & \texttt{vendored\_fetch} 离线宏包 23 件（\texttt{a90978ab}） & 绝版宏包缺件面 492 格失败$\to$73.2\% 纯净，纯净 +101 格 \\
09-17 & \texttt{illegal\_unit} 段修波（\texttt{d2c377bd}） & 108/108 闭环零回归，纯净率 84.76$\to$85.03\% \\
09-17 & ReDoS 原子组修复（\texttt{746237fd}） & held6 卡死 30min+$\to$0.1s；pytest 全量 48min$\to$3:46 \\
09-17 & 非锚定 [n] 译文解析撤除（\texttt{164a9e0f}，P0） & 译文静默错配通道封死，同批核销 4 个 P1 \\
09-17 & realn200 真臂验收（\texttt{f2847c33}） & 联合出 PDF 197/200=98.5\%，与模拟臂打平 \\
09-18 & c2-dollar 车道收口 & 泄漏 57$\to$0/136{,}049 译段（核心层口径单类清零） \\
09-18 & \texttt{gate\_scorecard} schema v3 & loop1 复跑 97.89/84.99 分毫不差——评测器自证稳定 \\
09-18 & qualbase esa2（$n$=1200） & 质量基线建档 均分 92.9 [92.4,93.3]，争议率 8.1\% \\
09-18 & loop2 冻结快照（$n$=5219） & 计分卡 97.26\%/86.42\%（纯净率 +1.43pt） \\
09-19 & corpus 8 层扩层（\texttt{13c603cf}） & 语料 5232$\to$13{,}266，EVAL\_ONLY 治理落地 \\
09-19 & 全量综合测试（本报告） & 12 臂同日落地；基线地板升至 90.2\%；stagerun 98.75\%/88.75\% \\
\bottomrule
\end{tabular}
\end{table}
      % 指标演化史（时间线图 + 归因表 + 大事记）
\section{2026-09-19 全量画像}

本节为 corpus\_v3 扩充收官（13{,}266 篇）后的首轮全量综合测试，12 个评测臂同日落地。总表见表~\ref{tab:snapshot}，逐节展开关键分布。

\begin{table}[htbp]
\centering\scriptsize
\caption{2026-09-19 综合测试逐臂指标总表}
\label{tab:snapshot}
\setlength{\tabcolsep}{4pt}
\begin{tabular}{@{}ll>{\raggedright\arraybackslash}p{6.4cm}>{\raggedright\arraybackslash}p{3.0cm}@{}}
\toprule
臂 & 规模 & 核心指标 & 门槛对比 \\
\midrule
完整性审计 & 13{,}266 格 & sha256 13{,}266/13{,}266，0 重复、0 缺失；13 个 stub 为 B07 有意构造 & EVAL\_ONLY 生效 \\
parsebench & 13{,}253 篇 / 28{,}904 .tex & 解析成功 \textbf{100.0\%}；一致率 \textbf{99.99\%}；泄漏 \textbf{0.004\%}（76/1{,}845{,}338）；expand 层足迹 strict 28105/norm 259/div 539；合平 93.7\% & 一致率/泄漏过门；\textbf{死占位符=3 未过}；\textbf{合平未过} \\
compilebench 裸编 & 500 格 $\times$2 & xelatex 纯净 56.3\%/出 PDF 80.3\%；tectonic 38.6\%/62.4\%；联合出 PDF \textbf{90.4\%} & --- \\
compilebench 中文链 & 500 格 & xelatex \textbf{72.4\%/85.4\%}（纯净 +16pt）；tectonic 45.7\%/59.8\% & --- \\
修复循环（v2 40） & 80 格 & xelatex 34 格失败$\to$\textbf{25 救回 74\%}；tectonic 1/16（EPS 路由主导） & --- \\
alignbench & 812 对 & 门全过；765/812 对无具名锚点，真实覆盖 42 对（保留率中位 1.0） & 覆盖不足 \\
e2e\_mock & corpus39 & 管线 xel \textbf{33/39 纯净} vs 裸编 xel 19/39；引入回归 3 & --- \\
e2e\_real & n=24+60 & 58/58 翻译；译段通过 6053/7580；回填残留 \textbf{0}；管线 xel 纯净 43/失败 7；修复循环再救 13；\textbf{0 引入回归} & --- \\
xlatbench & 271 样本 & 硬契约 \textbf{93.7\%}；http 错误 0；延迟 p50 4.2s/p95 29.7s & --- \\
qualbench & 324 篇/1937 译段 & 评审均分 \textbf{94.0}，中位 95；$\geq$90 占 90.6\%；争议率 5.5\%；critical 1 & --- \\
gullet & 200 文档 & 199/200 展开成功；中位 39.6ms/19 步 & --- \\
stagerun 流水线 & 80 篇跨层 & \textbf{联合出 PDF 79/80=98.75\%}；纯净 71/80=88.75\% & M2 出 PDF 门\textbf{过} \\
\bottomrule
\end{tabular}
\end{table}

\subsection{全链漏斗：stagerun 五段流水线}

五段批量流水线对 80 篇跨层样本的实测漏斗（图~\ref{fig:funnel}）：获取 80 $\to$ 解析 79 通过（1 拒）$\to$ 翻译模拟臂 76 通过 + 3 带瑕 $\to$ 编译中文臂 61 纯净 + 13 带瑕 + 5 失败 + 1 跳过 $\to$ 修复循环在 21 格上动作（失败$\to$纯净 5、带瑕$\to$纯净 5、带瑕保持 8、纯净复核 3）。联合终态 \textbf{纯净 71 + 带瑕 8 = 出 PDF 79/80（98.75\%）}，过 M2 联合出 PDF $\geq90\%$ 门；纯净率 88.75\% 距 90\% 门差 1.25pt。同格对照臂裸编译（50 纯净 + 19 带瑕 + 10 失败）给出归因锚点：管线加修复循环相对裸编译 \textbf{出 PDF +10 格（69$\to$79）、纯净 +21 格（50$\to$71）}。

\begin{figure}[htbp]
\centering
\begin{tikzpicture}
\begin{groupplot}[
  group style={group size=2 by 1,horizontal sep=10mm},
]
\nextgroupplot[
  xbar,
  width=7.0cm, height=4.6cm,
  xmin=0, xmax=95,
  xlabel={格数}, ytick=data,
  yticklabels={联合出 PDF,编译中文臂,翻译通过,解析通过,获取},
  y tick label style={font=\scriptsize},
  nodes near coords, nodes near coords style={font=\scriptsize},
  bar width=4.2mm,
  x label style={font=\small},
]
\addplot+[fill=cA!70,draw=cA] coordinates {(79,0) (74,1) (76,2) (79,3) (80,4)};
\nextgroupplot[
  ybar, bar width=5mm,
  width=5.4cm, height=4.6cm,
  ymin=0, ymax=90,
  symbolic x coords={裸编译,中文链+修复},
  xtick=data, x tick label style={font=\scriptsize},
  ylabel={格数}, y label style={font=\small},
  legend style={at={(0.5,-0.26)},anchor=north,legend columns=2,font=\scriptsize},
  nodes near coords, nodes near coords style={font=\scriptsize},
]
\addplot+[fill=cD!70,draw=cD] coordinates {(裸编译,50) (中文链+修复,71)};
\addplot+[fill=cC!75,draw=cC] coordinates {(裸编译,19) (中文链+修复,8)};
\legend{纯净,带瑕}
\end{groupplot}
\end{tikzpicture}
\caption{stagerun 五段漏斗与同格对照（80 篇跨层样本，翻译上游为模拟臂）。左：主流水线中文臂逐段存活数；右：同一 80 格上裸编译 vs 中文链+修复循环联合终态——管线相对裸编译出 PDF +10 格、纯净 +21 格（修复循环把 5 格失败全救成纯净、5 格带瑕升纯净）。}
\label{fig:funnel}
\end{figure}

\subsection{翻译质量：qualbench ESA 评分}

swe-2-medium 译文对 324 篇论文 1{,}937 个译段的 LLM 评审评分（初裁 swe-2-max，复核 swe-2-high，争议率 5.5\% 触发复核）：均分 94.0 / 中位 95，$\geq$90 分占 90.6\%，critical 级缺陷仅 1 例（漏译类）。分数段与文类分布见图~\ref{fig:qual}：节标题最高（96.9），图题 93.2，正文段 92.2——长散文段仍是质量洼地。

\begin{figure}[htbp]
\centering
\begin{tikzpicture}
\begin{groupplot}[
  group style={group size=2 by 1,horizontal sep=14mm},
  width=7.2cm, height=5cm,
]
\nextgroupplot[
  ybar, bar width=7mm,
  ymin=0, enlarge y limits={upper,value=0.15},
  symbolic x coords={$\geq$90,75--89,55--74,$<$55},
  xtick=data, x tick label style={font=\scriptsize},
  ylabel={译段数}, y label style={font=\small},
  nodes near coords, nodes near coords style={font=\scriptsize},
]
\addplot+[fill=cA!70,draw=cA] coordinates {($\geq$90,1756) (75--89,133) (55--74,41) ($<$55,7)};
\nextgroupplot[
  ybar, bar width=5mm,
  ymin=88, ymax=98,
  symbolic x coords={正文段,图题,节标题},
  xtick=data, x tick label style={font=\scriptsize,rotate=15},
  ylabel={均分}, y label style={font=\small},
  nodes near coords, nodes near coords style={font=\scriptsize},
]
\addplot+[fill=cD!70,draw=cD] coordinates {(正文段,92.2) (图题,93.2) (节标题,96.9)};
\end{groupplot}
\end{tikzpicture}
\caption{qualbench 评分分布（左：分数段；右：分类均分）。1{,}756/1{,}937 译段 $\geq$90；正文段是主要拉低项（缺陷榜前四：术语不一致 229、语法 197、误译 132、幻觉内容 66）。}
\label{fig:qual}
\end{figure}

\subsection{编译臂：裸编译 vs 中文链条件}

500 格抽样双引擎对照（图~\ref{fig:compile}）：中文链条件（产品链规范化 + ctex 注入后编译）反而比裸编译高 16pt 纯净率——规范化顺带修复了源文件的编码/格式缺陷。tectonic 在中文臂纯净率升而出 PDF 略降，EPS 图源是其死穴（见 \S\ref{sec:failures}）。

\begin{figure}[htbp]
\centering
\begin{tikzpicture}
\begin{axis}[
  ybar, bar width=6mm,
  width=11.5cm, height=5.4cm,
  ymin=0, ymax=100,
  enlarge x limits=0.35,
  ylabel={\%}, y label style={font=\small},
  symbolic x coords={xelatex 纯净,xelatex 出PDF,tectonic 纯净,tectonic 出PDF,联合出PDF},
  xtick=data, x tick label style={font=\scriptsize,rotate=28,anchor=east},
  legend style={at={(0.5,-0.36)},anchor=north,legend columns=2,font=\scriptsize},
  nodes near coords, nodes near coords style={font=\scriptsize},
]
\addplot+[fill=black!45,draw=black!60] coordinates {(xelatex 纯净,56.3) (xelatex 出PDF,80.3) (tectonic 纯净,38.6) (tectonic 出PDF,62.4) (联合出PDF,90.4)};
\addplot+[fill=cA!70,draw=cA] coordinates {(xelatex 纯净,72.4) (xelatex 出PDF,85.4) (tectonic 纯净,45.7) (tectonic 出PDF,59.8) (联合出PDF,0)};
\legend{裸编译基线, 中文链（规范化+ctex 后）}
\end{axis}
\end{tikzpicture}
\caption{compilebench 500 格：裸编译 vs 中文链条件。中文臂 xelatex 纯净率 +16.1pt——规范化顺带修复了源级缺陷；联合出 PDF 90.4\%（双引擎取优）。}
\label{fig:compile}
\end{figure}

\subsection{翻译硬契约：xlatbench}

271 个抽样上下文对 swe-2-medium 的硬契约测试：硬契约通过 254/271=93.7\%，http 错误 0。延迟 p50 4.2s / p90 16.2s / p95 29.7s / max 104.4s——批量档入闸排队在长尾可见。token 经济：输入 77.4k / 输出 40.8k（平均每译段 285/150），推理量均 93 tok，提示词缓存命中 0\%（抽样源分散）。按类契约率：压力样例 4/4、正文段 64/67、图题 62/67、条目 61/66、脚注 58/60——脚注/条目类是短板。失败类：\texttt{cs\_dropped}$\times$8、\texttt{ph\_miss}$\times$5、\texttt{validator}$\times$3、\texttt{ord\_hard}$\times$1；\texttt{ord\_soft} 78 条软换序告警（不破契约，但提示换序倾向）。

\subsection{端到端全链：模拟与真实双臂}

模拟臂（corpus39）：管线 xel 33/39 纯净 vs 裸编 xel 19/39——管线净正，翻译+回填步骤不损反增（规范化顺带修复）；3 格引入回归。真实臂 n=60：58/58 翻译完成，译段通过 6053/7580（带瑕 43/故障 8），回填残留占位符 0（过门），管线 xel 纯净 43/失败 7/带瑕 6，修复循环再救 13 格，\textbf{0 管线引入回归}。编译失败类：缺文件$\times$7、缺图$\times$3、信号终止$\times$3、引擎崩溃$\times$3——以语料源缺资产为主。
     % 2026-09-19 全量画像（逐臂 + 漏斗 + 分布）
\section{失败归因图谱}
\label{sec:failures}

\subsection{解析侧失败面}

全量 28{,}904 文件仅 1 个解析错误（语料源怪癖：伪装成 \texttt{.tex} 的 tar 包）与 1 个一致率分叉文件——硬失败面已枯竭，残面转移为告警谱（不挡解析但值得跟踪的形态）：\texttt{stray\_end} 3{,}007、\texttt{missing\_input} 2{,}757、\texttt{unpaired\_dollar} 383、\texttt{expand\_tail\_dropped} 313、\texttt{unclosed\_env} 223、\texttt{def\_parse\_fail} 140、\texttt{cite\_range\_key} 62、\texttt{env\_mismatch} 27 等。泄漏 76 条命中构成为 dollar 族 75 条 + 引用族 1 条——dollar 族长尾未穷举（演化史见图~\ref{fig:parse}）。

\subsection{编译侧失败类}

真实臂 n=60 的编译失败归因（图~\ref{fig:failclass} 左）：\texttt{missing\_file}$\times$7 与 \texttt{missing\_graphic}$\times$3 合计过半——语料源自身缺资产（arXiv 源包打包不全），非管线引入；\texttt{killed\_by\_signal}$\times$3 为超时/资源类；\texttt{driver\_fatal}$\times$3 为引擎级崩溃。

compilebench 500 格裸编译首错归因（修复循环 F 类命名）：xelatex 侧 \textbf{\texttt{missing\_file}$\times$94 居首}（代表 \texttt{psfig.sty}——古早宏包缺件）$>$ \texttt{other}$\times$48 $>$ \texttt{driver\_fatal}$\times$13；tectonic 侧 \textbf{\texttt{ps\_image}$\times$142 一家独大}（EPS 图源）。路由收益有实证：「eps/pstricks$\to$xel 优先」带 162 篇，xel 出 PDF 90/带瑕 24/失败 47 vs tec 4/22/135；\texttt{latex209\_suspect} 带 65 篇呈反向互补（xel 失败 43 vs tec 带瑕 36——古文在 tectonic 下更易出带瑕疵 PDF），\texttt{non\_utf8} 带 37 篇。学科面难度分层明显：cs 最易（xel 出 PDF 83/86）、math 次之（79/87）、astro-ph 最难（41/64）。

修复循环救回矩阵（corpus\_v2 40 篇对照）：xelatex 34 个失败格救回 25（74\%），主力规则为 \texttt{install\_file}（14 次触发，8 格出 PDF）、\texttt{install\_tfm}（12 次触发，7 格）、\texttt{static\_precheck}（40 格全触发）；装包榜首位 \texttt{subfigure.sty}$\times$10、\texttt{revtex4.cls}$\times$9。tectonic 臂仅救 1/16——\textbf{EPS 路由拒编是结构性死穴}（11 格全拒）：EPS 图源在 tectonic/xdvipdfmx 通路无解，唯一出路是路由到 xelatex。另有 \texttt{latex209\_reject} 闸拦下 LaTeX2.09 古文 8 格（该类论文需 209$\to$2e 重写通道，非装包可救）。本批 stagerun 中修复循环动作谱更轻量：21 格触发标准三件套全命中（\texttt{tar\_blob\_extract}、\texttt{subfile\_docclass\_strip}、\texttt{static\_precheck}），外加 \texttt{missing\_char\_fix}、\texttt{vendored\_fetch} 等 10 条各 1--4 次。

\subsection{翻译侧失败类}

xlatbench 失败类分布（图~\ref{fig:failclass} 右，17 个硬失败全归因）：\textbf{\texttt{cs\_dropped}$\times$8 居首}——控制序列在译文中丢失，且高度集中为脆弱间距命令（\verb|\ |、\verb|\,|、\verb|~| 被中文熔合吞掉，占 8 中 7 条）；\texttt{ph\_miss}$\times$5（占位符未还原，其中一格连丢 \texttt{CMD\_77}--\texttt{CMD\_85} 九个占位符——单译段级联性丢失而非均匀散布）、\texttt{comment\_eof}$\times$2（译文尾段注释未终结吞掉下一行字面量，可致回填后失控）、\texttt{validator}$\times$3、\texttt{ord\_hard}$\times$1；另 \texttt{ord\_soft} 78 条为不破契约的译段顺序软告警（合法中文换序的常见形态）。

qualbench 缺陷榜（1{,}937 已评译段，按严重度分层）：术语不一致 229、语法 197、误译 132、流畅度语域 99、幻觉内容 66、漏译片段 37、过度翻译 23、占位符破损 8。按族×严重度看结构：术语族 263 条（231 轻症/32 重症）是最大缺陷类但多为轻症；\textbf{约定族「不可译」25 条中 23 条重症}——不该译的内容被译（参考文献/命令名），是少量但系统性的协议违例；critical 级仅漏译 1 例。

\begin{figure}[htbp]
\centering
\begin{tikzpicture}
\begin{groupplot}[
  group style={group size=2 by 1,horizontal sep=16mm},
  width=7.4cm, height=5.2cm,
]
\nextgroupplot[
  xbar, bar width=3.6mm,
  xmin=0, enlarge x limits={upper,value=0.25},
  symbolic y coords={driver\_fatal,killed\_by\_signal,missing\_graphic,missing\_file},
  ytick=data, y tick label style={font=\scriptsize},
  xlabel={格数}, x label style={font=\small},
  nodes near coords, nodes near coords style={font=\scriptsize},
]
\addplot+[fill=cB!70,draw=cB] coordinates {(3,driver\_fatal) (3,killed\_by\_signal) (3,missing\_graphic) (7,missing\_file)};
\nextgroupplot[
  xbar, bar width=3.6mm,
  xmin=0, enlarge x limits={upper,value=0.15},
  symbolic y coords={ord\_soft,ord\_hard,validator,ph\_miss,cs\_dropped},
  ytick=data, y tick label style={font=\scriptsize},
  xlabel={样本数}, x label style={font=\small},
  nodes near coords, nodes near coords style={font=\scriptsize},
]
\addplot+[fill=cC!75,draw=cC] coordinates {(78,ord\_soft) (1,ord\_hard) (3,validator) (5,ph\_miss) (8,cs\_dropped)};
\end{groupplot}
\end{tikzpicture}
\caption{失败类分布。左：真实臂 n=60 编译失败归因——语料源缺资产主导；右：xlatbench 271 样本翻译失败/告警类——\texttt{cs\_dropped} 为头号硬失败。}
\label{fig:failclass}
\end{figure}

\subsection{门槛未过项与已知缺口}

\begin{enumerate}
  \item \textbf{死占位符=3}（parsebench）：死占位符残留 3 例，门槛为 0——唯一未过门项，值小但需逐例归因。全量 28{,}904 文件中唯一解析错误为「伪装成 \texttt{.tex} 的 tar 包」（\texttt{2211.00110} 内 \texttt{affinity\_mat.tex}）——语料源怪癖而非解析器缺陷；一致率唯一分叉文件同待逐例归因。
  \item \textbf{合平 93.7\% $<99\%$}：已知口径事项——多层语料中附属 \texttt{.tex}（孤儿文件，不被主文档引用）不进合平面的占比，口径而非缺陷。
  \item \textbf{alignbench 覆盖薄}：812 对中 765 对具名锚点为零——保留率指标实际只在 42 个 hyperref 对上有意义，指标口径需复核（分母该是所有对还是仅 hyperref 对）。
  \item \textbf{tectonic 短板}：EPS 路由无解 + biblatex$\leftrightarrow$biber 版本错配——EPS/bib 论文在 tectonic 下是先天缺省，路由层应前置判走 xelatex。
  \item \textbf{\texttt{cs\_dropped} 归因待定}：提示词侧约束不足 vs 网关侧截断，需在 xlatbench 加对照实验分辨。
  \item \textbf{stagerun 计分卡口径}：落盘记录混存首轮欠采的过期跳过行会污染读数（本批 989 条把联合出 PDF 冲成 7.4\% 假象；按获取样本口径重算 98.75\%）——工具层待加 \texttt{--since}/样本过滤参数。
  \item \textbf{未测臂}：translators\_bench、validbench（B6）、nightwatch、rundiff、triage、stage\_timing、wave、status\_panel——监控/辅助工具为主，本轮未覆盖。
\end{enumerate}
     % 失败归因图谱 + 已知缺口
\section{结论要点}
\label{sec:takeaway}

把六天 1{,}648 个 commit 与 12 臂综合读数压成十条可复用结论。每条给出数据锚点与对后续工作的含义。

\begin{enumerate}
  \item \textbf{管线是净正的，不是成本项}。同一 80 格上，中文链+修复循环相对裸编译出 PDF +10 格、纯净 +21 格；compilebench 500 格中文臂纯净率反超裸编 +16.1pt。规范化+修复循环顺带修好了源文件自身的缺陷——「翻译引入的编译代价」被「规范化修复的源缺陷」覆盖后还有盈余。含义：对外口径可以直接说「过本管线的论文比原样编译更容易出 PDF」。

  \item \textbf{修复循环是单一最大杠杆，且收益可复现}。09-16 接线当日联合出 PDF +18.9pt（同口径复跑确认）；v2 对照 40 篇 34 个失败格救回 25（74\%）；本批真实臂 7 个管线失败格全数救回。143 条规则中装包/补件类贡献主要救回量。含义：规则沉淀速度（31$\to$143 条/5 天）就是出 PDF 率的爬坡速度，机制台账$\to$规则这条反馈环值得继续投入。

  \item \textbf{解析硬失败面已经枯竭，残面转为告警谱}。28{,}904 文件解析成功 100\%、一致率 99.99\%，唯一解析错误是语料源怪癖（tar 伪装成 .tex）。剩下的是告警级形态（\texttt{stray\_end}、\texttt{missing\_input} 各三千量级）与泄漏长尾（dollar 族 75 条）。含义：解析段从「修 bug」阶段进入「清长尾」阶段，优先级应让位于翻译与编译侧。

  \item \textbf{泄漏归零是样本框的函数，不是终态}。dollar 族在 1{,}955 文件核心层清零后，8 倍体量全量复测重现 76 条命中。教训：任何「清零」结论必须绑定样本框声明；核心层绿不代表长尾绿，全量复测节奏要保持。

  \item \textbf{翻译硬失败高度集中，提示词可定向修}。\texttt{cs\_dropped} 8 例中 7 例是 \verb|\ |、\verb|\,|、\verb|~| 三个脆弱间距命令被中文熔合吞掉——不是均匀散布的鲁棒性问题，是单一可点名机制。含义：在提示词里对这三个命令显式钉保护即可预期消灭该类大半；\texttt{ph\_miss} 的级联丢失（一格连丢 9 个占位符）提示要查译段级联失效而非逐点丢失。

  \item \textbf{质量缺陷呈结构性分布：术语是大头轻症、约定违例是小头重症}。术语族 263 条但 231 条轻症；「不可译」约定 25 条中 23 条重症——参考文献与命令名被误译是系统性协议违例。含义：术语表资产值得继续扩（轻症面大、边际收益高），约定违例则要靠校验规则硬拦（量小但每条都是产品级事故）。

  \item \textbf{「出 PDF $\neq$ 成功」有实证兜底}。Mode B 注入 132 处占位符破坏，编译前被校验链 132/132 全捕获。这是整条链路最关键的非对称保证：翻译错了宁可不出 PDF，不能出一份静默错位的 PDF。含义：校验层投入（检出率、误报率、延迟三指标全绿）是产品正确性的根基，不是开销。

  \item \textbf{评测器自证过稳定性}。计分卡 schema v3 重放 loop1 得 97.89/84.99 分毫不差；报告全链路可复算（落盘记录追加式 + 样本集过滤口径）。含义：跨日对比读数可信，「指标涨了」可归因到代码而非测量噪声——这是多日迭代能收敛的前提。

  \item \textbf{基线自己会动}。裸编译地板四天涨约 19pt（71.7$\to$90.2），原因是离线宏包资产与 tlmgr usertree 积累——工具链成熟本身抬高了所有臂的分母。含义：任何「相对基线的增益」都要标注基线时点；「裸引擎撑死七成」这类早期结论会过期，README 与文档里的静态断言要定期复测。

  \item \textbf{未达标项与短板是显式清单，不是模糊感受}。M2 纯净率差 1.25pt（88.75 vs 90）、死占位符 3 例、合平口径 93.7\%、alignbench 有效覆盖仅 42 对、tectonic 的 EPS/biber 双死穴、\texttt{cs\_dropped} 归因待定、计分卡过期行污染——七项各有明确下一步（归因、口径修正、路由前置、对照实验、工具参数）。含义：下一里程碑的验收条件已经写死，不存在「感觉差不多了」的空间。
\end{enumerate}
     % 十条结论
\appendix
\section{口径与数据说明}

\subsection{统计口径纪律}

\begin{itemize}
  \item 每个数字标注规模 $n$ 与样本框；不同 $n$/口径的读数不连成同一演化线（图注注明口径切换点）。
  \item 计分卡读数一律按获取样本集过滤后计算——stagerun 落盘记录为追加式，跨轮重跑会残留过期行，原始行计数会把未参与样本的跳过计入分母。
  \item 编译终态三级：纯净 = 无告警出 PDF；带瑕 = 带可接受瑕疵出 PDF（\texttt{acceptable}、\texttt{best\_effort}、\texttt{dirty} 三档）；联合口径出 PDF = 纯净 + 带瑕。
  \item 延迟为网关端到端秒；批量档入闸排队约 120s，拥塞期 p95 会被队列时间放大，读数与空闲档不可比。
  \item qualbench 评分由 LLM 评审（swe-2-max）产出，争议判例经第二评审（swe-2-high）仲裁；评审自身偏差未做人工标定，分数应读作相对分布而非绝对正确率。
\end{itemize}

\subsection{随档清单与原始产物索引}

本报告图表全部由随档 \texttt{data/} 原始数据生成，不依赖本目录以外的图片路径；但 \texttt{refs/} 仅随档上一期复盘的两张 SVG，引用源文档的时点快照未随迁——对应档案位置见下表与 README.md（「随档 \texttt{refs/} 全收录」的原始表述以 README 为准）。

\begin{center}\small
\begin{tabular}{@{}l>{\raggedright\arraybackslash}p{10.6cm}@{}}
\toprule
目录 & 内容 \\
\midrule
\texttt{data/} & 六条时间线 CSV（计分卡 / 解析 / 语料 / 规则测试 / 翻译质量 / 每日 commit）；\texttt{arms\_history.csv} 377 行评测流水（含各轮 run 目录名、样本量、头条读数）；\texttt{commits.json} 1{,}648 commit 全录；\texttt{milestones.md} D0--D5 大事记 \\
\midrule
\texttt{refs/} & 随档仅 \texttt{metrics-timeline.svg} + \texttt{assets-growth.svg}（上一期复盘两图）。其余源文档快照未随迁，对应档案：评测器规格（七臂定义/门槛）→ \texttt{docs/spec/benchmark.md}；语料规格（层构造/治理）→ \texttt{docs/spec/corpus.md}；逐库横评协议 → \texttt{docs/dev/bench-harness.md}；验证分层契约 L0--L3 → \texttt{bench/TIERS.md}；层语义/EVAL\_ONLY → \texttt{bench/corpus/MANIFEST*.md} 族；机制台账快照 → \texttt{bench/corpus/mechanisms.jsonl}；8 库横评（D0 定案证据）→ \texttt{bench/archive-2026-09-20/results/00-grand-comparison.md}；本批 Markdown 总表 → \texttt{docs/research/methods/2026-09-19-bench-metrics-corpusv3.md}；上一期复盘正文与车队台账快照 → 迁移前存档 \texttt{tmp/old-docs-2026-09-20/docs/research/metrics-2026-09-19/refs/} \\
\bottomrule
\end{tabular}
\end{center}

各臂原始评测产物（编译工作区、译文 jsonl、PDF 中间件）为大体量文件不入档；溯源时按 \texttt{data/arms\_history.csv} 的 run 目录名在评测产物存档区查找对应轮次。

\subsection{过程事故记录（评测工具链教训）}

本批测试记录的评测驱动层教训，供后续批次回避：xlatbench 的 \texttt{manifest.parent} 必须等于语料根（跨目录清单会致格解析失败）；\texttt{--per-kind} 是全局限额非分层配额；pgrep 监控须用 \texttt{[x]pattern} 括号形防自匹配；compilebench \texttt{--gen-sample} 为分段设计（采样后即返回，编译需无参二次调用）；e2e 同 \texttt{--tag} 冻结样本，重抽样须换 tag+seed；gullet \texttt{--out} 续跑会保留过期错误行、须换新目录；stagerun 下游阶段须显式传 \texttt{--ids} 防跨层欠采；多会话共仓提交须走私有 \texttt{GIT\_INDEX\_FILE}（共享索引有 sweep/reset 竞争窗）。
     % 口径说明 + 结果目录索引

\end{document}
